Este estudo examina em que condições dispositivos controlados pelo utilizador podem partilhar o processamento de um assistente de IA mantendo limites explícitos de localização dos dados e autorização. A avaliação compara execução local, gateway HTTP e delegação peer-to-peer cumulativa de etapas, com tarefas, artefactos de modelos, configurações e condições de hardware declarados. O retrato atual inclui controlos concluídos de pedidos digitados e uma cadeia real de áudio pré-gravado que executou reconhecimento de fala, inferência de linguagem e síntese de fala. A cadeia gerou áudio de saída válido, mas falhou a tarefa de referência porque a instrução reconhecida diferiu da entrada esperada. Esse resultado negativo é preservado separadamente dos controlos digitados bem-sucedidos. As medições completas das condições de áudio e as comparações reais entre dispositivos permanecem incompletas; não se afirma latência de cauda de áudio, aceleração causal, privacidade global de implantação ou benefício geral da delegação. A contribuição nesta etapa é uma arquitetura de partilha implementada e um método auditável que distingue correção funcional, localização da execução, viabilidade de recursos e latência. As conclusões finais dependem de completar as condições não-iOS declaradas e apresentar falhas e limitações juntamente com os sucessos.

\textbf{Palavras-chave:} Sistemas Distribuídos; Redes P2P; Localização dos Dados; Assistentes de Voz.
